\documentclass[12pt,a4paper]{article}

\usepackage[a4paper,margin=1.7cm]{geometry}
\usepackage[T2A]{fontenc}
\usepackage[utf8]{inputenc}
\usepackage[russian]{babel}
\usepackage{amsmath,amssymb,mathtools}
\usepackage{array,booktabs,longtable}
\usepackage{xcolor}
\usepackage{hyperref}
\usepackage{tikz}

\hypersetup{
  colorlinks=true,
  linkcolor=blue!55!black,
  urlcolor=blue!55!black,
  pdfborder={0 0 0}
}

\setlength{\parindent}{0pt}
\setlength{\parskip}{0.7em}
\setlength{\tabcolsep}{3pt}
\renewcommand{\arraystretch}{1.25}

\newcommand{\tasktitle}[1]{\section*{Задание №#1}}
\newcommand{\subhead}[1]{\par\medskip\textbf{#1}\par\smallskip}

\begin{document}

\begin{titlepage}
  \centering
  \vspace*{0.25\textheight}
  {\LARGE\bfseries Ревью домашней работы\par}
  \vspace{1.8cm}
  {\large Автор: Лёвин Артём Александрович\par}
  \vspace{0.8cm}
  {\large 4 октября 2026 года\par}
  \vfill
  \begin{tikzpicture}
    \draw[line width=0.7pt] (-3.8,0) .. controls (-1.7,0.65) and (1.7,-0.65) .. (3.8,0);
  \end{tikzpicture}
\end{titlepage}

\newpage
\section*{Сводная таблица}
\small
\begin{longtable}{>{\centering\arraybackslash}p{0.9cm} p{2.6cm} p{5.0cm} p{3.3cm} p{3.3cm}}
\toprule
\textbf{№} & \textbf{Тема} & \textbf{Суть ошибки} & \textbf{Ваш ответ} & \textbf{Правильный ответ} \\
\midrule
\endfirsthead
\toprule
\textbf{№} & \textbf{Тема} & \textbf{Суть ошибки} & \textbf{Ваш ответ} & \textbf{Правильный ответ} \\
\midrule
\endhead
\hyperref[task:2]{2} & Линейное тепловое расширение & Прибавку длины $4$ мм приняли за полную длину стержня при температуре $T$. & Не указан отдельной строкой. & $25^{\circ}\mathrm{C}$ \\
\midrule
\hyperref[task:4]{4} & Квадратичная зависимость пути от времени & Разность двух пройденных расстояний заменена выражением через квадрат разности времён. & $0{,}2$ м & $1{,}4$ м \\
\midrule
\hyperref[task:9]{9} & Объём цилиндра и высота уровня воды & Четверть первоначального объёма не была переведена в четверть первоначальной высоты воды. & Из решения получены $t=0$ и $t=800$. & $200$ с \\
\bottomrule
\end{longtable}
\normalsize

\newpage
\thispagestyle{empty}
\vspace*{\fill}
\begin{center}
  {\LARGE\bfseries Разбор ошибок}
\end{center}
\vspace*{\fill}

\newpage
\phantomsection\label{task:2}
\tasktitle{2}

\subhead{Текст задания}
Длина стержня при температуре $T$ задаётся формулой
\[
L(T)=L_0(1+\alpha T),\qquad L_0=8\text{ м},\qquad \alpha=2\cdot10^{-5}\,^{\circ}\mathrm{C}^{-1}.
\]
При какой температуре стержень будет на $4$ мм длиннее, чем при $0^{\circ}\mathrm{C}$?

\subhead{Ваше решение}
В решении правильно переведено $4$ мм в $0{,}004$ м, после чего записано
\[
0{,}004=8\bigl(1+2\cdot10^{-5}T\bigr).
\]
Далее из этого равенства выражается $T$ через отрицательное число.

\subhead{Ваш ответ}
Не указан отдельной строкой.

\subhead{Ошибка}
\textcolor{red}{Ошибка:} при составлении уравнения прибавка длины $0{,}004$ м была приравнена к полной длине стержня $L(T)$.

\subhead{Где именно возникает ошибка}
Последний корректный шаг --- перевод единиц:
\[
4\text{ мм}=0{,}004\text{ м}.
\]
Первый неверный шаг:
\[
0{,}004=8\bigl(1+2\cdot10^{-5}T\bigr).
\]
Правая часть по условию обозначает \emph{всю} длину стержня при температуре $T$, а число $0{,}004$ м обозначает только то, насколько эта длина должна увеличиться по сравнению с длиной при $0^{\circ}\mathrm{C}$.

\subhead{Почему это неверно}
При $0^{\circ}\mathrm{C}$ длина стержня равна
\[
L(0)=8\bigl(1+2\cdot10^{-5}\cdot0\bigr)=8\text{ м}.
\]
Фраза «на $4$ мм длиннее» означает, что к исходным $8$ м нужно прибавить $0{,}004$ м. Поэтому искомая полная длина равна $8{,}004$ м, а не $0{,}004$ м.

Иначе можно работать сразу с приростом длины. Из формулы
\[
L(T)=L_0+L_0\alpha T
\]
видно, что прирост равен $L_0\alpha T$. Именно его и надо приравнивать к $0{,}004$ м.

\subhead{Ключевая идея}
\textcolor{blue!60!black}{Ключевая идея:} сначала нужно различить исходную длину, полную новую длину и прибавку длины. Число $4$ мм относится к прибавке, а формула $L(T)$ даёт полную длину.

\subhead{Исправление от последнего правильного шага}
После перевода $4$ мм в метры можно записать
\[
8{,}004=8\bigl(1+2\cdot10^{-5}T\bigr).
\]
Вычтем исходные $8$ м:
\[
0{,}004=8\cdot2\cdot10^{-5}T.
\]
Тогда
\[
0{,}004=16\cdot10^{-5}T=0{,}00016T,
\]
откуда
\[
T=\frac{0{,}004}{0{,}00016}=25.
\]

\subhead{Полное правильное решение}
При нулевой температуре длина стержня равна $8$ м. Требуется, чтобы он стал длиннее на $4$ мм, то есть на $0{,}004$ м. Следовательно,
\[
L(T)=8{,}004\text{ м}.
\]
Подставим это значение в формулу:
\[
8{,}004=8\bigl(1+2\cdot10^{-5}T\bigr).
\]
Разделим обе части на $8$:
\[
1{,}0005=1+2\cdot10^{-5}T.
\]
Вычтем единицу:
\[
0{,}0005=2\cdot10^{-5}T.
\]
Отсюда
\[
T=\frac{0{,}0005}{2\cdot10^{-5}}=25.
\]
Значит, требуемая температура равна $25^{\circ}\mathrm{C}$.

\subhead{Проверка}
При $T=25$ получаем
\[
L(25)=8\bigl(1+2\cdot10^{-5}\cdot25\bigr)
=8(1+0{,}0005)=8\cdot1{,}0005=8{,}004\text{ м}.
\]
Разница с исходной длиной действительно равна $0{,}004$ м, то есть $4$ мм.

\subhead{Итог}
\textcolor{teal!60!black}{Итог:} правильный ответ --- $25^{\circ}\mathrm{C}$.

\subhead{Задача на закрепление}
Длина стержня задаётся формулой $L(T)=L_0(1+\alpha T)$, где $L_0=10$ м, $\alpha=1{,}5\cdot10^{-5}\,^{\circ}\mathrm{C}^{-1}$. При какой температуре стержень будет на $3$ мм длиннее, чем при $0^{\circ}\mathrm{C}$?

\newpage
\phantomsection\label{task:4}
\tasktitle{4}

\subhead{Текст задания}
Камень падает до поверхности воды; расстояние падения задаётся формулой $h(t)=5t^2$ м. До дождя падение длилось $0{,}8$ с, после дождя --- $0{,}6$ с. На сколько метров поднялся уровень воды?

\subhead{Ваше решение}
В решении записано
\[
\Delta h=5\Delta t^2=5\cdot0{,}2^2=0{,}2\text{ м}.
\]

\subhead{Ваш ответ}
$0{,}2$ м.

\subhead{Ошибка}
\textcolor{red}{Ошибка:} разность двух расстояний $h(0{,}8)-h(0{,}6)$ была заменена величиной $h(0{,}8-0{,}6)$.

\subhead{Где именно возникает ошибка}
Верно найдено, что время падения уменьшилось на
\[
0{,}8-0{,}6=0{,}2\text{ с}.
\]
Первый неверный шаг --- запись
\[
\Delta h=5\Delta t^2.
\]
Для квадратичной зависимости это не работает: разность значений функции не равна значению функции от разности аргументов.

\subhead{Почему это неверно}
До дождя поверхность воды находилась ниже, поэтому камень падал $0{,}8$ с. Расстояние от точки падения камня до старого уровня воды равно
\[
h(0{,}8)=5\cdot0{,}8^2.
\]
После дождя уровень воды поднялся, и расстояние стало
\[
h(0{,}6)=5\cdot0{,}6^2.
\]
Подъём уровня воды --- это разность именно этих двух расстояний:
\[
\Delta h=h(0{,}8)-h(0{,}6).
\]

Алгебраически причина та же:
\[
0{,}8^2-0{,}6^2\ne(0{,}8-0{,}6)^2.
\]
Левая часть равна $0{,}64-0{,}36=0{,}28$, а правая --- только $0{,}04$.

\subhead{Ключевая идея}
\textcolor{blue!60!black}{Ключевая идея:} когда величина зависит от времени квадратично, сначала вычисляют значение величины для каждого момента времени, а уже затем находят разность полученных значений.

\subhead{Исправление от последнего правильного шага}
Используем оба времени отдельно:
\[
\Delta h=5\cdot0{,}8^2-5\cdot0{,}6^2.
\]
Вычислим:
\[
5\cdot0{,}8^2=5\cdot0{,}64=3{,}2,
\]
\[
5\cdot0{,}6^2=5\cdot0{,}36=1{,}8.
\]
Поэтому
\[
\Delta h=3{,}2-1{,}8=1{,}4\text{ м}.
\]

\subhead{Полное правильное решение}
До дождя камень за $0{,}8$ с проходил расстояние
\[
h_1=5\cdot0{,}8^2=3{,}2\text{ м}.
\]
После дождя время падения стало $0{,}6$ с, значит новое расстояние до воды равно
\[
h_2=5\cdot0{,}6^2=1{,}8\text{ м}.
\]
Поверхность воды стала ближе к точке начала падения на
\[
h_1-h_2=3{,}2-1{,}8=1{,}4\text{ м}.
\]
Следовательно, уровень воды поднялся на $1{,}4$ м.

\subhead{Проверка}
После дождя время падения уменьшилось, поэтому расстояние до поверхности воды должно уменьшиться. Полученные значения $3{,}2$ м и $1{,}8$ м этому соответствуют, а их разность положительна и равна $1{,}4$ м.

\subhead{Итог}
\textcolor{teal!60!black}{Итог:} уровень воды поднялся на $1{,}4$ м.

\subhead{Задача на закрепление}
Камень падает до поверхности воды, причём расстояние падения задаётся формулой $h(t)=5t^2$ м. До подъёма уровня воды падение длилось $1$ с, а после подъёма --- $0{,}8$ с. На сколько метров поднялся уровень воды?

\newpage
\phantomsection\label{task:9}
\tasktitle{9}

\subhead{Текст задания}
Для вертикального цилиндрического бака высота воды при сливе задаётся формулой
\[
h(t)=16\left(1-\frac{t}{400}\right)^2,\qquad 0\le t\le400,
\]
где $h$ --- в метрах, $t$ --- в секундах. Через сколько секунд после открытия крана останется четверть первоначального объёма воды?

\subhead{Ваше решение}
В решении записано
\[
16=16\left(1-\frac{t}{400}\right)^2,
\]
после чего получено
\[
\left(1-\frac{t}{400}\right)^2=1
\]
и рассмотрены значения $t=0$ и $t=800$.

\subhead{Ваш ответ}
Из решения получены значения $t=0$ и $t=800$; отдельный окончательный ответ не выделен.

\subhead{Ошибка}
\textcolor{red}{Ошибка:} условие о четверти первоначального объёма было заменено условием о сохранении первоначальной высоты воды.

\subhead{Где именно возникает ошибка}
Первый неверный шаг:
\[
16=16\left(1-\frac{t}{400}\right)^2.
\]
Число $16$ м --- это первоначальная высота воды. Если в цилиндрическом баке остаётся четверть первоначального объёма, то при неизменной площади основания остаётся и четверть первоначальной высоты, то есть $4$ м.

\subhead{Почему это неверно}
Объём воды в вертикальном цилиндре равен произведению площади основания $S$ на высоту воды $h$:
\[
V=Sh.
\]
Площадь основания не меняется. Поэтому отношение объёмов равно отношению высот:
\[
\frac{V(t)}{V(0)}=\frac{h(t)}{h(0)}.
\]
Если объём стал равен одной четверти первоначального, то
\[
h(t)=\frac14h(0)=\frac14\cdot16=4\text{ м}.
\]

Запись с $16$ слева фактически означает $h(t)=h(0)$, то есть бак всё ещё содержит первоначальный объём. Именно поэтому одно из полученных значений --- $t=0$, момент начала слива. Второе значение $t=800$ вообще выходит за заданный промежуток $0\le t\le400$.

\subhead{Ключевая идея}
\textcolor{blue!60!black}{Ключевая идея:} в цилиндрическом баке объём прямо пропорционален высоте воды, потому что площадь основания постоянна. Поэтому четверть объёма соответствует четверти высоты.

\subhead{Исправление от последнего правильного шага}
Первоначальная высота равна
\[
h(0)=16\text{ м}.
\]
Четверть этой высоты:
\[
\frac{16}{4}=4\text{ м}.
\]
Поэтому нужно решать уравнение
\[
4=16\left(1-\frac{t}{400}\right)^2.
\]
Разделим на $16$:
\[
\left(1-\frac{t}{400}\right)^2=\frac14.
\]
На промежутке $0\le t\le400$ выражение $1-\frac{t}{400}$ неотрицательно, поэтому
\[
1-\frac{t}{400}=\frac12.
\]
Тогда
\[
\frac{t}{400}=\frac12,
\qquad t=200.
\]

\subhead{Полное правильное решение}
Пусть площадь основания цилиндрического бака равна $S$. Тогда первоначальный объём воды равен
\[
V_0=16S.
\]
Четверть первоначального объёма равна
\[
\frac{V_0}{4}=4S.
\]
Так как текущий объём равен $V(t)=Sh(t)$, условие $V(t)=4S$ равносильно условию
\[
h(t)=4.
\]
Подставим формулу высоты:
\[
4=16\left(1-\frac{t}{400}\right)^2.
\]
Получаем
\[
\left(1-\frac{t}{400}\right)^2=\frac14.
\]
Поскольку $0\le t\le400$, имеем
\[
0\le1-\frac{t}{400}\le1.
\]
Следовательно, берём неотрицательное значение:
\[
1-\frac{t}{400}=\frac12.
\]
Отсюда
\[
t=200\text{ с}.
\]

\subhead{Проверка}
Подставим $t=200$:
\[
h(200)=16\left(1-\frac{200}{400}\right)^2
=16\left(\frac12\right)^2
=4\text{ м}.
\]
Высота стала в четыре раза меньше первоначальной. Для цилиндра это означает, что и объём стал в четыре раза меньше, то есть осталась ровно четверть первоначального объёма.

\subhead{Итог}
\textcolor{teal!60!black}{Итог:} четверть первоначального объёма останется через $200$ секунд.

\subhead{Задача на закрепление}
В вертикальном цилиндрическом баке высота воды задаётся формулой
\[
h(t)=25\left(1-\frac{t}{500}\right)^2,\qquad 0\le t\le500.
\]
Через сколько секунд после начала слива в баке останется четверть первоначального объёма воды?

\vfill
\begin{center}
\small Лёвин Артём Александрович эксклюзивно для Григория Архипова
\end{center}

% Краткое сообщение Григорию по итогам ревью:
% Григорий, в работе есть три момента, которые важно отметить: во втором задании прибавка длины была принята за полную длину стержня; в четвёртом разность двух пройденных расстояний заменена квадратом разности времени; в девятом четверть объёма воды не была переведена в четверть высоты.

\end{document}
